\section{A geometric criterion for the two endpoint moments}
Consider the repeated bases and fixed block-count bound of Section 17.
Keep positivity, normalization and the uniform stationary $L^{15/7}$ bound.
Centering is now optional. Along any subsequence, pass to a fixed number of
blocks with each size either constant or tending to infinity. Let $J$ index
the growing blocks. For each nonempty $J$ consider every weak limit $\lambda$
of the stationary vector of standardized sums on $J$. Such vectors are tight
by H\"older and the subgaussian uniform block laws.

\paragraph{Criterion.} Each family $(Ph)^{-15}$ and $(P\sqrt\pi)^{-15}$
is uniformly integrable if and only if every such $\lambda$ assigns
positive mass to $\{w:u\cdot w>0\}$ for every unit $u$ on $J$.
If one of the limits is supported in a closed half-space through zero,
the expectations for both families diverge along a further subsequence.

For sufficiency, a failure of uniform integrability yields a subsequence
with tail thresholds tending to infinity and tail integrals bounded below.
Extract block sizes and $\lambda$ as above. For each $u$ choose a bounded
open subset of $u\cdot w>2\epsilon_u$ with positive $\lambda$-mass.
Directions in a neighborhood of $u$ have projection above $\epsilon_u$ on
that subset. A finite sphere cover and the open-set bound for weak convergence
give fixed $\epsilon,c,M>0$ along the extracted sequence: every unit direction
has stationary mass $\ge c$ in its positive slice of the radius-$M$ ball,
with projection $\ge\epsilon$. Section 17 transfers this mass to both
inputs $h,\sqrt\pi$. Apply (BX4) in growing blocks and the uniform positive
row lower bound in bounded blocks. The shell sum after (BX5) contradicts the
chosen tails. With $J$ empty use the entire space and the bounded-block
lower bounds.

For necessity suppose $\lambda(u\cdot w\le0)=1$. For $f=\sqrt\pi$ or
$f=h$ and any set $A$,
\[
                 \left(\int_A f\,d\mu\right)^2\le\nu(A).
\]
For $h$ this follows from
$\langle h1_A,Ph\rangle\ge\langle h1_A,P(h1_A)\rangle
\ge(\int_A h\,d\mu)^2$, using entrywise positivity and the nonnegative
spectrum with constant mode. For $\sqrt\pi$ use Cauchy--Schwarz.
Thus the projected measures $f\,d\mu$ are tight subprobabilities.
Extract weak limits $\eta_f$ simultaneously. Each is supported in the same
half-space: exhaust its open complement by bounded open sets whose closures
are inside it and apply the mass bound and weak convergence on the closures.

Write $\rho=1/4$, $q=1-\rho^2$. In one growing block the row-to-uniform
layer ratio for a compact input $z$ is bounded uniformly in the output $w$.
Indeed reversibility expresses it as the reversed row probability at $z$
divided by its uniform layer mass, at least $c_R/\sqrt n$ for $|z|\le R$.
The reversed row is a convolution with a binomial of size $\ge n/2$ and
fixed success probability $(1\pm\rho)/2$, so every atom is $O(n^{-1/2})$.
On compact sets of both variables the ratio converges uniformly on the
lattice to $q^{-1/2}\exp[-\rho^2(z^2+w^2)/(2q)+\rho zw/q]$.
For details, in plus-count coordinates each Bernoulli characteristic factor
has absolute square $1-q\sin^2(t/2)$. After centering, Fourier inversion at
scale $1/\sqrt n$ has an integrable Gaussian majorant on its whole
rescaled period. Its logarithm on bounded frequencies equals
$-qt^2/8+O(|t|^3/\sqrt n)$ in the rescaled variable. The variance in
plus counts is $nq/4$, and the mean in standardized sign units is $\rho z$.
The same calculation for uniform counts has $q=1$ and zero sign mean.
Dividing the local estimates yields the stated kernel.

Multiply over $J$ and bound all fixed-block factors by their maximum.
For bounded convergent inputs $z_n\to z$ in the growing blocks, truncating
the output to a compact set, then using tightness and the global row bound
for the remaining tail, gives
\[
 \limsup Pf(z_n)\le C\int q^{-|J|/2}
 e^{-\rho^2(|z|^2+|w|^2)/(2q)+\rho z\cdot w/q}\,d\eta_f(w)
\]
uniformly in bounded-block inputs. Here $Pf(z_n)$ denotes the layer value
with those additional inputs. If $z=tu+y$, $t\ge0$, $y\perp u$, $|y|\le1$,
the cross term with $tu$ is nonpositive on the support. Completing the square
in $w$ bounds the right side by $C'e^{-\rho^2|z|^2/(2q)}$, independently
of the cylinder length. Couple the uniform growing block sums by independent
quantiles to a standard Gaussian vector. On every bounded portion of this
cylinder, Fatou's inequality bounds the liminf of the inverse moment below
by a constant times its Euclidean volume: the Gaussian density cancels the
quadratic exponent because $15\rho^2=q$. Letting the length increase proves
divergence, also in dimension one. The positive transverse width is essential
in higher dimensions.
